\section{Teoría Programación Dinámica}

% primer frase pero perdió contra la del Goat del rey leon
% ``Aquellos que no recuerdan el pasado están condenados a repetirlo'' - George Santayana (1905)
% por qué guardamos datos en caché/memoización.

``El pasado puede doler, pero tal como yo lo veo, puedes huir de él o aprender de él''. - Rafiki (El Rey León, 1994).

\textbf{Programación Dinámica} es una técnica de programación (como los algoritmos Greedy). Ayuda a dejar de repetir cálculos aprovechando los ya hechos, usar resultados pasados para resolver el futuro. 
La solución por Programación Dinámica brinda resultados, estos resultados calculados los podemos guardar, lo que nos puede ayudar a optimizar el tiempo de ejecución. Esa técnica de guardado se llama \textit{Memoization} o \textit{Memorización}.

La Programación Dinámica explora implícitamente el espacio de posibles soluciones, lo cual siempre va a convenir ver el problema general en pequeños \textit{subproblemas} del mismo.
Nos basamos en la intuición que nos da la Memoization para reconocer los subproblemas y utilizarlos para construir la solución. Una vez que tenemos todas las soluciones memorizadas, el problema está resuelto. Nos ayuda evitar explorar un espacio exponencial de soluciones por fuerza bruta, y de esa manera podemos reducir la complejidad temporal del algoritmo a implementar.

Los \textit{subproblemas} establecen una lógica inductiva:

\noindent- ``Si yo puedo asumir que ya tengo calculados todas las soluciones anteriores (casos más pequeños), ¿Cómo puedo usarlos para construir la que necesito?''

Esto nos lleva a pensar en un cierto \textit{patrón} que pueden generar los subproblemas y, en consecuencia, una \textit{Ecuación de Recurrencia}\footnote{Una ecuación de recurrencia en el ámbito de la informática es una expresión matemática que define el tiempo de ejecución o el costo de un algoritmo en función de sí mismo para entradas más pequeñas.} que cumpla con el armado de los subproblemas.

¿Qué pasaría si, buscando un valor anterior para encontrar la solución requerida, viene la persona más importante y que más sabe de este tema a decirte: ``Esta es la solución que necesitas.''? Confiás plenamente, pero ¿cómo utilizarías esta información a tu favor?

\noindent Ejemplo: 

Secuencia de Fibonacci. Para un número $n$, no es necesario calcular todas las ramas de sus anteriores, porque se van a \textit{repetir} resultados. Viene esta persona y te dice: ``Para los 2 casos que quieras sumar, ya calculé los resultados y son estos''. Entonces, simplemente hay que sumar esos dos \textit{óptimos} que nos brindó, que sabemos que son correctos y no dudamos en ningún momento.
¿Cuál es la forma de los subproblemas? ¿Cuál es la ecuación de recurrencia?
\begin{equation*} 
    \mathcal{OPT}(n) = \mathcal{OPT}(n-1) + \mathcal{OPT}(n-2)
\end{equation*}

Siendo $OPT(n)$ el resultado de fibonacci que queremos calcular. 

Si $n = 5$, calcula $fibonacci(5) = fibonacci(4) + fibonacci(3)$ 

Entonces, si \textit{iterativamente} tenemos precalculados, recordados, \textit{memorizados} todos los resultados que queremos averiguar, los fibonacci anteriores a $n$, en una simple suma obtenemos el resultado requerido. 

Para la resolución de los problemas utilizando la técnica de Programación Dinámica se prefiere mucho la forma iterativa (bottom-up), y no utilizar la forma recursiva (top-down). Similar a los algoritmos Greedy.
Esto es debido a que:
\begin{enumerate}
    \item Es mucho más fácil de entender qué esta pasando, cuándo y cómo.
    \item Como consecuencia de lo anterior, se facilita el cálculo de la complejidad (temporal y espacial).
    \item La resolución de problemas con esta técnica suelen ser más fáciles de resolver de forma inductiva que recursiva.
    \item Nos estructura mucho más cómo pensar la solución\footnote{Primero la ecuación de recurrencia, después la programación del problema (siendo este último paso, sencillo)}.
    \item Se pueden aplicar optimizaciones de memoria fácilmente.
\end{enumerate}

¿Como se va a saber que puedo utilizar Programación Dinámica en un problema?
\begin{itemize}
    \item    Hay un número polinomial de subproblemas
    \item La solución al problema original puede ser construido a partir de soluciones a subproblemas
    \item Hay un orden natural de los subproblemas de menor a mayor. Los subproblemas “mayores” son resueltos mediante la composición de problemas “menores”.
\end{itemize}

Otra utilidad que nos sirve usar la técnica de \textit{Memoization},comúnmente, permite pasar de un tiempo de ejecución exponencial a uno polinomial. Además, nos permite en ciertos ejercicios una optimización de memoria: En el anterior ejemplo de fibonacci, se puede utilizar un arreglo con todos los óptimos de $0$ hasta $n$, para recordar las soluciones. Pero mejor aún, utilizar una optimización y guardar solamente $2$ variables recordando las últimas dos soluciones. Eso nos va a brindar una complejidad espacial constante de $O(1)$.

Otra sección relacionada a la técnica de Programación Dinámica, es la \textit{reconstrucción del problema}. Si yo quiero recordar y devolver todos los resultados del fibonacci en el orden en que los usé, ahí sí se utiliza un arreglo con todos los óptimos precalculados y en una recorrida se pueden obtener qué resultados y \textit{cuándo} los utilicé. La reconstrucción del problema es construir la solución de abajo hacia arriba.

\subsection{Programación Dinámica aplicada al problema}

En el contexto del juego de las monedas, el desafío de Sophia es planificar la mejor estrategia. El futuro es planificable y predecible, ya que sabemos que Mateo agarra siempre la moneda de mayor valor que se le cruce en frente. 

Entonces, la decisión de Sophia en el turno actual no reduce el tamaño del problema en $1$ moneda, sino en $2$ monedas (cada ronda). Su propia elección sumada a la reacción avariciosa de Mateo.

\textbf{El Pensamiento de Sophia:} \\
En vez de elegir sin pensar la moneda de mayor valor (como haría un algoritmo Greedy), Sophia piensa mentalmente los dos universos posibles en su turno:
\begin{enumerate}
    \item \textbf{¿Qué pasa si elijo la moneda de la izquierda?} Mateo, por instinto greedy, va a agarrar la mayor de las dos monedas que queden en los extremos. Sabiendo exactamente cuál va a elegir Mateo, Sophia se pregunta: \textit{¿Cuántos puntos máximos me aseguro con las monedas que sobran en la mesa?}.
    \item \textbf{¿Qué pasa si elijo la moneda de la derecha?} Mateo va a hacer exactamente lo mismo. Sophia vuelve a calcular: \textit{En este otro caso, ¿cuántos puntos me aseguro para el futuro?}.
\end{enumerate}

Sophia simplemente calcula el total de cada camino (la moneda que elige ahora + los puntos que ya sabe que va a ganar en el futuro) y ejecuta la opción que le devuelva el número más grande.

\vspace{1ex}
\noindent \textbf{Ecuación de Recurrencia:} \\
Para formalizar este pensamiento, definimos a $\mathcal{O}\mathcal{P}\mathcal{T}(\text{prim}, \text{ult})$ como el máximo valor acumulado que Sophia puede asegurarse en la fila de monedas actual.

\textbf{Casos Base (y comúnmente, el final del juego):}
\begin{itemize}
    \item Si queda sólo una moneda: Sophia se la queda. $\mathcal{O}\mathcal{P}\mathcal{T}(\text{prim}, \text{prim}) = m[\text{prim}]$
    \item Si quedan dos monedas: Sophia elige la mayor. $\mathcal{O}\mathcal{P}\mathcal{T}(\text{prim}, \text{ult}) = \max(m[\text{prim}], m[\text{ult}])$
\end{itemize}

\textbf{Caso General (Cuando quedan 3 o más monedas):} \\
Lo mejor es simplemente el máximo entre esos dos futuros casos que calcula Sophia mentalmente:
\[
\mathcal{O}\mathcal{P}\mathcal{T}(\text{prim}, \text{ult}) = \max 
\begin{cases} 
    m[\text{prim}] + \text{Futuro si agarré la Primera} \\ 
    m[\text{ult}] + \text{Futuro si agarré la Última}
\end{cases}
\]

Matemáticamente, ese "futuro" depende del Greedy aplicado de Mateo, quedando la ecuación completa así:
\[
\mathcal{O}\mathcal{P}\mathcal{T}(\text{prim}, \text{ult}) = \max 
\begin{cases} 
    m[\text{prim}] + 
    \begin{cases} 
        \mathcal{O}\mathcal{P}\mathcal{T}(\text{prim} + 2, \text{ult}) & \text{si Mateo elige } m[\text{prim} + 1] \\ 
        \mathcal{O}\mathcal{P}\mathcal{T}(\text{prim} + 1, \text{ult} - 1) & \text{si Mateo elige } m[\text{ult}] 
    \end{cases} \\[1.5em]
    
    m[\text{ult}] + 
    \begin{cases} 
        \mathcal{O}\mathcal{P}\mathcal{T}(\text{prim} + 1, \text{ult} - 1) & \text{si Mateo elige } m[\text{prim}] \\ 
        \mathcal{O}\mathcal{P}\mathcal{T}(\text{prim}, \text{ult} - 2) & \text{si Mateo elige } m[\text{ult} - 1] 
    \end{cases} 
\end{cases}
\]

\noindent \textbf{Demostración de mayor valor acumulado posible}:

Para demostrar que esta estrategia es la mejor que le podemos recomendar a Sophia y garantiza siempre el mayor valor acumulado posible para Sophia, nos basamos en el razonamiento inductivo. 

Como vimos en los Casos Base, es un obvio que Sophia hace la decisión perfecta cuando quedan 1 o 2 monedas en la fila. El éxito de la estrategia de Programación Dinámica se basa en que, cuando Sophia se enfrenta a un problema de 3, 4 o $n$ monedas, ella no improvisa. Al evaluar el "camino izquierdo" y el "camino derecho", la ecuación consulta los resultados de los subproblemas más chicos (el futuro). Ventaja superior comparado con el pensamiento Greedy de agarrar siempre la de mayor valor.

Como esos subproblemas del futuro fueron resueltos previamente construidos desde los casos base hacia arriba, iterativamente, Sophia está sumando su moneda actual con un futuro que es matemáticamente insuperable. Dado que la ecuación evalúa el 100\% de las ramas posibles y siempre extrae la función $\max()$ entre ellas, es lógicamente imposible que exista una combinación de elecciones en la fila que sume un valor mayor al que la estrategia da.

Se va a verificar, por demostración matemática y con ejemplos, en una sección a continuación.

\subsection{Pensamiento y utilidades} % utilidades ¿? no me convence.

El Pensamiento de Programación Dinámica es análisis. La mente comúnmente tiende a abrumarse con el futuro, esta mentalidad introduce: entender que los problemas gigantes no son bloques fijos, sino un árbol de subproblemas vinculados por el tiempo. 
Adoptar esta forma de razonar implica aprender a pensar de atrás hacia adelante. En lugar de pararse en el casillero inicial y avanzar a ciegas tanteando opciones a ver que onda, el ``pensador dinámico'' se para conceptualmente en la meta final que quiera y reconstruye el camino de forma inversa. Y lo mejor de todo, las personas hacen esto todo el tiempo en el día a día cuando planifican su vida, partiendo sus objetivos grandes en pequeñas acciones concretas hasta llegar a la primera decisión que tiene al alcance para hacer hoy. (Aristóteles hablando sobre la deliberación en Ética nicomáquea\footnote{Expone en el Libro III de esta obra que la deliberación no empieza por el principio, sino por la meta. Analizamos los requisitos de fin a principio, descubriendo que ``el último paso en la resolución es el primero en la ejecución''.})

Esta forma de razonar retroespectiva tiene una base analítica muy fuerte: si el futuro está planificado de manera óptima, el presente se resuelve solo. Trasladado a las decisiones de la vida real, esto nos da una estructura clara para manejar la optimización del esfuerzo del día a día y la satisfacción aplazada a través de ``ecuaciones de recurrencia personales''. El que adopta este chip mental entiende que las decisiones actuales son variables que modifican el escenario del mañana. Por eso, se acepta de forma natural un costo, un esfuerzo extra o un rendimiento subóptimo en el plano inmediato (como aguantarse de salidas por estudiar o meterle horas a un proyecto) con la certeza matemática de que es el único camino válido para maximizar el resultado global a largo plazo. (Principio de Optimalidad de Bellman)

En un plano cognitivo, el enfoque PD introduce la \textit{memoization} como una técnica estricta para guardar la experiencia personal. El pensamiento dinámico no deja el desperdicio del aprendizaje ni el desgaste de tropezar dos veces con la misma piedra. Si o si exige que cada crisis superada, cada error corregido y cada subproblema resuelto quede anotado en una suerte de ``tabla mental'' permanente. Al cruzarse con situaciones repetitivas en la rutina diaria, la mente no reinventa la lógica desde cero, sino que hace una consulta rápida a su memoria. Y así, la toma de decisiones cotidianas no permiten el azar y pasa a pasa a ser algo más razonable, donde lo vivido se transforma de forma directa en un mapa predictivo para resolver la incertidumbre.
